\section{公的年金等控除}

\subsection{公的年金等に係る雑所得の計算}

\begin{fillblank}[年金1 2018 問題1(3)\textcircled{4}（改題）]
次に定めるB さんについて、公的年金等に係る雑所得は、（　　A　　）万円となる。

なお、令和2年分年以降の公的年金等控除額を使用するものとする。

＜B さんの情報＞
\begin{itemize}
  \item 年齢：66 歳
  \item 公的年金等の収入金額：公的年金（年 70 万円）および通算企業年金（年 30 万円）
  \item 通算企業年金に関し、企業年金連合会に移換前の制度において加入者負担掛金はない。
\end{itemize}
\end{fillblank}

\begin{answerbox}
  A: 0
\end{answerbox}

\begin{fillblank}[年金1 2018 問題1(3)\textcircled{5}（改題）]
  次に定めるC さんについて、公的年金等に係る雑所得は、（　　A　　）万円となる。

  なお、令和2年分年以降の公的年金等控除額を使用するものとする。

＜C さんの情報＞
\begin{itemize}
  \item ・年齢：62 歳
\item 公的年金等の収入金額：確定給付企業年金の老齢給付金（年 130 万円）のみ
\item 確定給付企業年金において加入者負担掛金はない。
\end{itemize}

\end{fillblank}

\begin{answerbox}
  A: $130万円 \times 0.75 - 27.5万円 = 70万円$
\end{answerbox}

\begin{fillblank}[自作 2026/10/25]
  次に定めるD さんについて、公的年金等に係る雑所得は、（　　A　　）万円となる。

  なお、令和2年分年以降の公的年金等控除額を使用するものとする。

＜D さんの情報＞
\begin{itemize}
  \item ・年齢：62 歳
\item 公的年金等の収入金額：確定給付企業年金の老齢給付金（年 130 万円）のみ
\item 確定給付企業年金の老齢給付金（年 130 万円）のうち、加入者拠出分は年30万円とする。
\end{itemize}

\end{fillblank}

\begin{answerbox}
\begin{flalign*}
130万円 - 30万円 &= 100万円  \\
          100万円 \times 1.0 - 60万円 = 40万円 \\
  \end{flalign*}
    
\end{answerbox}

\begin{fillblank}[年金1 2019 問題1(4)\textcircled{5}（改題）]
  1961 年 10 月 1 日生まれの者が 2026 年の公的年金等の年間収入金額が 100 万円であったと
き、この者の2026年分の公的年金等に係る雑所得は、 （　　A　　）円となる。

  なお、令和2年分年以降の公的年金等控除額を使用するものとする。
\end{fillblank}

\begin{answerbox}
  2026年12月31日時点で65歳なので、65歳以降の速算表を用いると、公的年金等に係る雑所得は0円。
\end{answerbox}

\newpage

\subsection{その他}

\begin{truefalse}[年金1 H28 問題1(6)\textcircled{1}]
    \textbf{下線部の正誤を判定し、誤っている場合は正しい内容に改めなさい。}
\\  

  公的年金等に係る雑所得は公的年金等控除額を控除することができるため、他の雑所得より
も税制上優遇されている。また、公的年金等に係る雑所得の課税方法は\underline{分離課税}となっている。
\end{truefalse}

\begin{answerbox}
  \textbf{誤}

  総合課税
\end{answerbox}

\begin{explanationbox}
  退職所得は分離課税。公的年金等に係る雑所得は総合課税。
\end{explanationbox}

\begin{fillblank}[年金1 2022 問題1(6)(ア)]
    \noindent \textbf{◯所得税法}
  \begin{lawstructure}
    \begin{lawarticle}{第三十五条}{雑所得}
      雑所得とは、利子所得、配当所得、不動産所得、事業所得、給与所得、退職所得、山林所得、譲渡所得及び一時所得のいずれにも該当しない所得をいう。
      
      \lawparagraph{2}
      雑所得の金額は、次の各号に掲げる金額の合計額とする。

      \lawitem{一}
      その年中の公的年金等の収入金額から公的年金等控除額を控除した残額

      \lawitem{二}
      その年中の雑所得（公的年金等に係るものを除く。）に係る総収入金額から必要経費を控除した金額

      \lawparagraph{3}
      （略）

      \lawparagraph{4}
      第二項に規定する公的年金等控除額は、次の各号に掲げる場合の区分に応じ当該各号に定める金額とする。

      \lawitem{一}
      その年中の公的年金等の収入金額がないものとして計算した場合における第二条第一項第三十号（定義）に規定する合計所得金額（次号及び第三号において「公的年金等に係る雑所得以外の合計所得金額」という。）が千万円以下である場合　次に掲げる金額の合計額（当該合計額が六十万円に満たない場合には、六十万円）

      \hangedlawitem{イ}{3}
      （　　A　　）万円
      
      \hangedlawitem{ロ}{3}
      その年中の公的年金等の収入金額から（　　B　　）万円を控除した残額の次に掲げる場合の区分に応じそれぞれ次に定める金額
    
      \lawitem{二、三　（略）}
    \end{lawarticle}
  \end{lawstructure}
\end{fillblank}

\begin{answerbox}
  \begin{enumerate}[A:]
    \item 四十
    \item 五十
  \end{enumerate}
\end{answerbox}

\begin{truefalse}[年金1 2018 問題2(3)\textcircled{4}]
       \textbf{以下の文章の正誤を判定し、誤っている場合は正しい内容に改めなさい。}
\\  

  平成30 年度税制改正における公的年金等控除の見直しにおいて、公的年金等に係る雑所
得以外の所得に係る合計所得金額が 1,000 万円を超える場合の控除額について、上限を設け
ることとされた。
\end{truefalse}

\begin{answerbox}
 公的年金等に係る雑所得以外の所得ではなく、公的年金等の収入金額が 1,000 万円を超える
場合の控除額について、上限が設けられることとされた。
\end{answerbox}

\begin{truefalse}[年金1 H29 問題2(3)\textcircled{3}]
         \textbf{以下の文章の正誤を判定し、誤っている場合は正しい内容に改めなさい。}
\\  

  平成 23 年分以後の各年分について、 公的年金等の収入金額が 400 万円以下で、かつ、公的年
金等に係る雑所得以外の所得金額が 20 万円以下である場合には、所得税の確定申告は不要とさ
れている。
\end{truefalse}

\begin{answerbox}
  \textbf{正}
\end{answerbox}